\chapter{Conclusions and Future Work}
\label{ch:conclusions}

This thesis has investigated the accuracy of opponent state estimation in autonomous racing, addressing the uncertainty associated with sensor measurements, the response of the tracking filter to abrupt changes in the opponent's dynamics, and the possibility of improving state reconstruction through non-causal estimation. The proposed methods have been developed and evaluated using experimental data collected during racing sessions, within the existing EKF-based target tracking architecture of the Dallara AV-24 and EAV-25 autonomous racecars.

A first significant outcome of this work is the experimental characterization of measurement errors under the actual operating conditions of the vehicles. The analysis has considered the errors produced by the complete perception pipeline, providing a physically meaningful and representative description of the measurement uncertainty encountered by the tracker. After accounting for reference and calibration errors, temporal and spectral analyses revealed that the remaining uncertainty cannot, in general, be adequately represented by a stationary white Gaussian noise model. Instead, several detection sources exhibit a composite error structure, consisting of a slowly varying, temporally correlated component superimposed on a faster, approximately decorrelated contribution. The autocorrelation of these errors was characterized using a first-order Gauss--Markov model for the slow component, together with an additional white-noise contribution. To investigate whether the slowly varying behavior could be explained by the measurement geometry, the relationship between the errors and the relative position and orientation of the opponent was subsequently examined. Although some dependencies were identified, they were not sufficiently reproducible across racing sessions to support a general deterministic compensation.

Building on this characterization, an innovation-based covariance matching approach has been developed to adapt the measurement noise covariance of the LiDAR detection pipelines online. Its experimental validation has demonstrated that the adaptive covariance provides an accurate approximation of the fast, decorrelated component of the measurement errors, offering a more representative description of its variability than a fixed covariance. This contribution remains superimposed on the slowly varying error component, which cannot be independently observed with the currently available sensor configuration.

The second part of the thesis has addressed the degradation of the estimated speed and acceleration during abrupt changes in the opponent's dynamics. A feed-forward adaptation of the process noise covariance has been developed to increase the responsiveness of the causal filter when a maneuver is detected, while retaining the nominal filtering behavior during steady motion. The experimental results have confirmed the effectiveness of this strategy during correctly identified maneuvers, while also highlighting the influence of the delay and noise affecting the jerk estimate used for detection. These effects introduce a trade-off between timely activation and spurious detections, which ultimately constrains the improvement achievable through a purely reactive approach.

A more substantial improvement in state reconstruction has been obtained through the implementation of a fixed-lag Rauch--Tung--Striebel smoother. By incorporating measurements acquired after the instant being estimated, the smoother refines the estimated motion and improves the temporal reconstruction of the opponent's maneuvers. With a smoothing window of \SI{300}{ms}, the speed RMSE is reduced by approximately \SI{31}{\percent} when all detection sources are available and by \SI{59}{\percent} when RaDAR measurements are excluded, relative to the corresponding causal estimates. The latter result is particularly significant, as it demonstrates the effectiveness of smoothing in reconstructing the target speed even when it must be inferred from successive position measurements, rather than directly observed through range-rate information. These improvements are achieved at the cost of a bounded estimation lag, with the refined estimate referring to a past instant rather than to the current state.

The most significant quantitative improvement has been achieved by combining the feed-forward process noise adaptation with the RTS smoother. In the evaluated scenarios, the speed RMSE decreases from \SI{1.069}{m/s} for the causal feed-forward estimate to \SI{0.356}{m/s} for the combined configuration, corresponding to a reduction of \SI{66.7}{\percent}. Similarly, the acceleration RMSE decreases from \SI{3.589}{m/s^2} to \SI{1.786}{m/s^2}, an improvement of \SI{50.2}{\percent}. Furthermore, compared with the RTS smoother applied to the nominal filter, the combined approach provides an additional \SI{43.8}{\percent} reduction in speed RMSE, from \SI{0.634}{m/s} to \SI{0.356}{m/s}. This result demonstrates the complementary nature of the two approaches: improving the response of the forward filter remains beneficial even when its estimates are subsequently refined using future measurements.

Overall, the results demonstrate that substantial improvements in target tracking accuracy can be achieved without replacing the underlying estimation architecture. The experimental characterization of sensor errors provides a more representative description of measurement uncertainty, while adaptive process noise and fixed-lag smoothing address the challenges posed by rapidly changing target dynamics. In particular, their combination enables a considerably more accurate reconstruction of the opponent's motion, showing how the complementary use of causal adaptation and non-causal estimation can improve tracking performance in demanding autonomous racing conditions.

\section{Future Work}
\label{sec:future_work}

A first extension concerns the application of innovation-based covariance matching to the RaDAR and camera detection sources. Their innovation sequences are affected by the slowly varying difference between their measurement biases and the LiDAR-related component already reflected in the target state estimate. A possible development is therefore to estimate and compensate these relative biases online, for instance by augmenting the filter with sensor-specific bias states modeled as first-order Gauss--Markov processes. Once the relative slow component has been compensated, covariance matching could be applied to the remaining fast, approximately decorrelated contribution. Integrating this procedure into the target tracking module would allow the measurement covariances of RaDAR and camera to adapt to their operating conditions, improving their relative weighting in the sensor fusion process and, in particular, enabling more effective use of the camera detections.

A second extension concerns the estimation of jerk for the feed-forward process noise adaptation. In the current implementation, jerk is obtained by differentiating the estimated acceleration, with additional filtering required to limit noise amplification. This introduces a trade-off between the delay of the maneuver detector and the occurrence of spurious activations. An alternative would be to adopt a motion model that explicitly includes jerk as a state variable, allowing this quantity to be estimated recursively rather than reconstructed by differentiation. The resulting jerk estimate could drive the same hysteresis-based process noise adaptation developed in this thesis. A comparison with the current approach, in terms of maneuver detection delay, false activations and speed and acceleration estimation errors, would determine whether the extended model provides a more reliable activation signal while preserving the nominal filter behavior outside maneuvers.
